\documentclass{article}

\usepackage{neurips_2026}

\usepackage[utf8]{inputenc}
\usepackage[T1]{fontenc}
\usepackage{hyperref}
\usepackage{url}
\usepackage{booktabs}
\usepackage{amsfonts}
\usepackage{amssymb}
\usepackage{amsmath}
\usepackage{nicefrac}
\usepackage{microtype}
\usepackage{xcolor}
\usepackage{algorithm}
\usepackage{algorithmic}
\usepackage{graphicx}
\usepackage{subcaption}
\usepackage{mathtools}
\usepackage{amsthm}

% Theorem environments
\newtheorem{theorem}{Theorem}
\newtheorem{lemma}[theorem]{Lemma}
\newtheorem{proposition}[theorem]{Proposition}
\newtheorem{corollary}[theorem]{Corollary}
\newtheorem{definition}[theorem]{Definition}
\newtheorem{remark}[theorem]{Remark}

% Shortcuts
\newcommand{\R}{\mathbb{R}}
\newcommand{\E}{\mathbb{E}}
\newcommand{\cG}{\mathcal{G}}
\newcommand{\cN}{\mathcal{N}}
\newcommand{\cM}{\mathcal{M}}
\newcommand{\bW}{\mathbf{W}}
\newcommand{\bx}{\mathbf{x}}
\newcommand{\lam}{\lambda_2}
% \newcommand{\phigr}{\varphi}  % removed: no longer needed for QRS-DR

\title{Deterministic Near-Ramanujan Topologies for Decentralized Learning}

\author{%
  Anonymous Authors
}

\begin{document}

\maketitle

% ============================================================================
% ABSTRACT
% ============================================================================
\begin{abstract}
Communication topology is a key determinant of convergence speed in decentralized
learning, yet practitioners either rely on random graphs---which offer no
worst-case guarantees---or hand-designed topologies that sacrifice spectral
quality at finite scale.  We introduce \emph{QRS-DR} (Quadratic Residue
Scatter with Degree Regularization), a fully deterministic algorithm that
builds $d$-regular graphs on $n$ nodes with algebraic connectivity $\lam$
near the Ramanujan bound $d - 2\sqrt{d-1}$ for any $n \ge 4$ and $d \ge 3$
with $nd$ even.  The scatter phase runs in $O(nd)$ time; degree
regularization adds $O(n^2 d)$ post-processing to guarantee exact
$d$-regularity.
%
The resulting $d$-regular graph can be decomposed into edge-disjoint
matchings via standard edge coloring, yielding a one-peer-per-round
communication schedule.
%
We prove convergence bounds for decentralized SGD on QRS-DR topologies
and show that at practically relevant scales ($n \le 1000$, $d \le 8$),
QRS-DR achieves larger spectral gaps than random $d$-regular graphs in
90\% of tested configurations while providing deterministic,
coordination-free, per-node-computable topology construction.
%
Experiments on CIFAR-10 and CIFAR-100 under heterogeneous data partitions
confirm that the improved spectral properties translate to faster
convergence and higher final accuracy in decentralized training.
%
Beyond spectral analysis, we conduct systems-level evaluation via SimGrid
simulation, demonstrating that QRS-DR topologies degrade gracefully under
heterogeneous link bandwidths, straggler nodes, and link failures---regimes
where finite-time-consensus topologies lose their theoretical guarantees.
\end{abstract}

% ============================================================================
% 1. INTRODUCTION  (~1.5 pages)
% ============================================================================
\section{Introduction}
\label{sec:intro}

% --- Hook: DFL is important, topology matters ---
Decentralized learning removes the central parameter server bottleneck by
having each worker communicate only with a fixed set of neighbors
\cite{lian2017dpsgd, koloskova2020unified}.
The communication topology---encoded as a mixing matrix $\bW$---directly
controls the convergence rate: the spectral gap $1 - |\lambda_2(\bW)|$
appears in every known convergence bound
\cite{koloskova2020unified, neglia2020topology}.

% --- Problem: existing topologies are inadequate ---
In practice, topologies are chosen ad hoc.
\emph{Ring} and \emph{torus} graphs are easy to construct but have spectral
gaps that vanish as $O(1/n^2)$.
\emph{Random $d$-regular graphs} achieve near-optimal spectral gaps with high
probability \cite{bordenave2015}, but offer no worst-case guarantees: a
single unlucky draw can produce a poorly connected graph, and reproducing the
same topology across all nodes requires a shared random seed and coordination
protocol.
\emph{Algebraic constructions} (LPS \cite{lps1988}, Margulis--Gabber--Galil
\cite{margulis1988}) achieve Ramanujan-optimal gaps but require $n$ to be
prime or a prime power, restricting practical deployment.

% --- Gap and contribution ---
Recent work on the Base-$(k{+}1)$ Graph \cite{takezawa2023base} achieves
\emph{finite-time exact consensus} via time-varying topologies, but
assumes homogeneous links and loses this guarantee under any communication
failure---a limitation explicitly acknowledged by the authors.
We ask: \emph{can we build a deterministic, near-Ramanujan $d$-regular
topology for arbitrary $(n,d)$, with a natural decomposition into
per-round communication schedules, that remains robust under realistic
network conditions?}

\paragraph{Contributions.}
\begin{enumerate}
\item \textbf{QRS-DR construction.}
    A fully deterministic algorithm that produces a $d$-regular graph on
    $n$ nodes with provably $\lam = \Omega(d)$ (constant spectral gap)
    and empirically $\lam \ge d - 2\sqrt{d-1}$ (near-Ramanujan) for any
    valid $(n,d)$ with $nd$ even, $n \ge 4$, $d \ge 3$.  Phase~1 (Quadratic
    Residue Scatter) uses degree-2 polynomial maps over $\mathbb{F}_p$
    to generate near-uniform edge candidates in $O(nd)$ time; Phase~2
    (Degree Regularization) deterministically adjusts the graph to exact
    $d$-regularity in $O(n^2 d)$ time with canonical tie-breaking
    (Section~\ref{sec:construction}).

\item \textbf{Matching decomposition for one-peer-per-round protocols.}
    The $d$-regular output graph can be decomposed into $d$ edge-disjoint
    perfect matchings via standard edge coloring (Vizing's theorem).
    Activating one matching per round yields a bandwidth-optimal
    communication schedule (Section~\ref{sec:matching}).

\item \textbf{Convergence analysis.}
    We prove $O\!\left(\frac{1}{\sqrt{nT}} +
    \frac{1}{T \cdot \gamma^2}\right)$ convergence for decentralized SGD
    on the QRS-DR topology, where $\gamma$ is the spectral gap
    of the gossip mixing matrix (Section~\ref{sec:theory}).

\item \textbf{Empirical validation.}
    On CIFAR-10/100 with heterogeneous data partitions (Dirichlet
    $\alpha \in \{0.1, 0.3, 1.0\}$), QRS-DR topologies
    converge faster and reach higher accuracy than ring, torus,
    exponential graph, Base-$(k{+}1)$ \cite{takezawa2023base}, and
    random $d$-regular baselines, with the largest gains in the
    high-heterogeneity regime (Section~\ref{sec:experiments}).

\item \textbf{Systems-level robustness via SimGrid simulation.}
    Using the SimGrid network simulator \cite{simgrid2014}, we evaluate
    wall-clock convergence under heterogeneous link bandwidths, straggler
    nodes, and random link failures (Section~\ref{sec:simgrid}).
    QRS-DR degrades gracefully in all regimes, while the
    Base-$(k{+}1)$ Graph---which achieves finite-time exact consensus
    under ideal conditions---loses its convergence guarantee when even a
    single link fails.
\end{enumerate}


% ============================================================================
% 2. RELATED WORK  (~1 page)
% ============================================================================
\section{Related Work}
\label{sec:related}

\paragraph{Decentralized SGD and topology design.}
Decentralized parallel SGD (D-PSGD) \cite{lian2017dpsgd} demonstrated
that peer-to-peer gradient averaging can match or outperform centralized
training on bandwidth-limited networks.  Subsequent work extended this
to asynchronous settings \cite{lian2018adpsgd} and directed graphs via
stochastic gradient push \cite{assran2019sgp}.
Koloskova et al.~\cite{koloskova2020unified} provided a unified
convergence framework showing that the spectral gap of the mixing
matrix directly controls the heterogeneity-dependent convergence term,
while Neglia et al.~\cite{neglia2020topology} confirmed empirically
that sparse topologies can accelerate convergence even without
communication delays.
MATCHA~\cite{wang2019matcha} decomposes the topology into matchings
and activates connectivity-critical links more frequently.
The \emph{exponential graph}~\cite{ying2021expgraph} achieves
$\tilde{O}(n^3)$ transient iterations with $O(\log n)$-degree static
or $O(1)$-degree one-peer schedules.
EquiTopo~\cite{jin2022equitopo} attains network-size-independent
$O(1)$ consensus rate with $\Theta(\ln n)$ degree, and
DSGD-CECA~\cite{ding2023ceca} achieves unit per-iteration
communication cost with $\tilde{O}(n^3)$ transient iterations for
arbitrary $n$.
The BlueFog framework~\cite{ying2021bluefog} provides a practical
PyTorch implementation supporting both static and dynamic topologies.
Most recently, Wang et al.~\cite{wang2025promise} identified three
key factors for practical speedups over AllReduce and demonstrated
wall-clock gains with decentralized Adam on up to 64 GPUs.

\paragraph{Spectral analysis, generalization, and data heterogeneity.}
Vogels et al.~\cite{vogels2022beyondgap} showed that the spectral gap
alone is not fully predictive of empirical performance in deep learning:
collaboration enables larger learning rates, and parameters like
effective neighborhood size also matter.
Zhu et al.~\cite{zhu2022topaware} provided the first topology-aware
generalization bound for D-SGD, proving that generalizability
correlates positively with spectral gap.
Le Bars et al.~\cite{lebars2023neighborhood} introduced
\emph{neighborhood heterogeneity} as the key quantity coupling topology
and data heterogeneity, and proposed learning data-dependent sparse
topologies.
On the algorithmic side, RelaySum~\cite{vogels2021relaysum} distributes
information uniformly via relay routing, making convergence depend on
diameter rather than spectral gap.
Quasi-global momentum~\cite{lin2021qgm} and
SlowMo~\cite{wang2020slowmo} improve performance under heterogeneity
via momentum-based corrections, while
D\textsuperscript{2}~\cite{tang2018d2} applies variance reduction to
decentralized training.
Epidemic Learning~\cite{devos2023epidemic} randomizes communication
partners each round, achieving $O(n^3/s^2)$ transient iterations when
each node contacts $s$ random peers.

\paragraph{Time-varying topologies and finite-time consensus.}
Takezawa et al.~\cite{takezawa2023base} introduced the Base-$(k{+}1)$
Graph, a time-varying topology achieving \emph{finite-time exact
consensus} in $O(\log_{k+1} n)$ rounds for any $n$ and maximum degree
$k$.  This is a powerful property---after a finite number of rounds,
all nodes hold exactly the average of their initial values---but as
acknowledged by the authors, the construction assumes homogeneous links
and \emph{loses its finite-time guarantee under any communication
failure}.
Nguyen et al.~\cite{nguyen2025gtft} formalized the connection between
finite-time consensus sequences and gradient tracking (GT-FT), showing
that convergence becomes independent of both function heterogeneity and
individual graph connectivity.  Fauth et al.~\cite{fauth2025approxftc}
extended this to \emph{approximate} finite-time consensus, showing that
even inexact sequences yield benefits.
Our work takes a complementary approach: rather than requiring exact
consensus in finite time (fragile to failures), QRS-DR
provides a \emph{static} topology with near-optimal spectral gap that
degrades gracefully under realistic network conditions.

\paragraph{Communication compression.}
Orthogonal to topology design, communication compression reduces
per-message cost.
Choco-SGD~\cite{koloskova2019choco} was the first gossip algorithm
supporting arbitrary compression with linear convergence;
Koloskova et al.~\cite{koloskova2020chocosgd} extended this to
non-convex deep learning at $O(1/\sqrt{nT})$.
DeepSqueeze~\cite{tang2019deepsqueeze} combined error compensation
with decentralization, and
PowerGossip~\cite{vogels2020powergossip} applied low-rank compression
to decentralized settings.
Our topology construction is compatible with these methods: the
per-round matching schedule transmits exactly one message per node,
making it straightforward to apply any compressor on top.

\paragraph{Expander graphs and Ramanujan bounds.}
The Ramanujan bound $\lam \ge d - 2\sqrt{d-1}$ is the optimal spectral
gap for infinite families of $d$-regular graphs, established by the
Alon--Boppana theorem~\cite{nilli1991, alon1986}.
The seminal constructions of Lubotzky--Phillips--Sarnak~\cite{lps1988}
and Margulis~\cite{margulis1988} produce Ramanujan graphs as Cayley
graphs of PGL$(2)$, but require the degree to be one plus a prime
(or prime power~\cite{morgenstern1994}).
Marcus, Spielman, and Srivastava~\cite{mss2015} proved existence of
bipartite Ramanujan graphs for all degrees via interlacing polynomials,
but their proof is non-constructive.
Friedman~\cite{friedman2008} proved that random $d$-regular graphs are
near-Ramanujan with high probability (Alon's conjecture);
Bordenave~\cite{bordenave2015} simplified and extended this to random
lifts.
The first \emph{explicit deterministic} near-Ramanujan construction for
all constant degrees is due to Mohanty, O'Donnell, and
Paredes~\cite{mohanty2020}, who derandomized Bordenave's random-lift
approach in poly$(n)$ time.
In the HPC context, SpectralFly~\cite{young2022spectralfly} showed that
LPS-based Ramanujan topologies outperform DragonFly and SlimFly
interconnects on spectral gap and bisection bandwidth.
Our QRS-DR construction differs from all the above: it is
deterministic, works for \emph{arbitrary} $(n, d)$ with $nd$ even
(not just special degrees or vertex counts), uses quadratic residue
arithmetic over $\mathbb{F}_p$ to achieve near-Ramanujan spectral gaps
with a theoretical anchor in Paley graph theory, and is per-node
computable---each node can reconstruct the full topology from $(n,d)$
alone without coordination.

\paragraph{Decentralized training at scale.}
Recent work has pushed decentralized training to foundation-model
scale.  DiLoCo~\cite{douillard2024diloco} uses federated averaging
with large inner steps to train language models with $500\times$ less
communication.  OpenDiLoCo~\cite{jaghouar2024opendiloco} and the
INTELLECT series~\cite{prime2024intellect1} scaled this to 10B+
parameters across continents.
SWARM Parallelism~\cite{ryabinin2023swarm} enables model-parallel
training on heterogeneous unreliable devices.
DecentLaM~\cite{yuan2022decentlam} addresses momentum bias in
large-batch decentralized training.
While these works focus on algorithmic innovations at scale, they
largely use ad-hoc or fixed topologies; our work provides the
graph-theoretic foundation for choosing optimal communication
topologies in such systems.

\paragraph{Systems-level evaluation.}
Most topology papers evaluate convergence in terms of communication
rounds, ignoring the physical network layer.  Recent work has begun
addressing this gap: bandwidth-aware topology
optimization~\cite{bandwidthaware2025} formulates topology design under
heterogeneous bandwidth constraints, and Marfoq et
al.~\cite{marfoq2020throughput} designed throughput-optimal topologies
for cross-silo federated learning.
SimGrid~\cite{simgrid2014} is a widely-used framework for simulating
distributed systems with realistic latency, bandwidth, and failure
models.  To our knowledge, we are the first to use SimGrid to evaluate
decentralized learning topologies under heterogeneous and fault-prone
network conditions, bridging the gap between graph-theoretic analysis
and systems-level performance.


% ============================================================================
% 3. PRELIMINARIES  (~0.5 pages)
% ============================================================================
\section{Preliminaries}
\label{sec:prelim}

\paragraph{Notation.}
We consider a network of $n$ workers connected by an undirected $d$-regular
graph $G = (V, E)$ with $|V| = n$.  The \emph{Laplacian} is
$L = D - A$ where $A$ is the adjacency matrix and $D = dI$.
Its eigenvalues are $0 = \lambda_1 \le \lambda_2 \le \cdots \le \lambda_n$.
The \emph{algebraic connectivity} $\lam(G)$ measures how well-connected
$G$ is; larger $\lam$ implies faster mixing.

\paragraph{Ramanujan bound.}
A $d$-regular graph is \emph{Ramanujan} if every eigenvalue $\lambda$ of its
adjacency matrix satisfies $|\lambda| \le 2\sqrt{d-1}$ (excluding $\pm d$).
Equivalently, $\lam(L) \ge d - 2\sqrt{d-1}$.  This is the optimal spectral
gap achievable by infinite families of $d$-regular graphs
\cite{lps1988, nilli1991}.

\paragraph{Decentralized SGD.}
Each worker $i$ holds a local loss $f_i(\bx)$ and maintains a local
parameter $\bx_i^{(t)}$.  At each round $t$:
\begin{equation}
\label{eq:dpsgd}
\bx_i^{(t+1)} = \sum_{j=1}^{n} W_{ij}^{(t)} \bx_j^{(t)}
              - \eta \, g_i^{(t)},
\end{equation}
where $g_i^{(t)}$ is a stochastic gradient and $\bW^{(t)}$ is a doubly
stochastic mixing matrix derived from the topology active at round $t$.


% ============================================================================
% 4. THE QRS-DR CONSTRUCTION  (~2 pages)
% ============================================================================
\section{The QRS-DR Construction}
\label{sec:construction}

The QRS-DR (Quadratic Residue Scatter + Degree Regularization)
construction builds a $d$-regular graph on $n$ nodes in two phases:
a \emph{scatter} phase that generates near-uniform edge candidates using
quadratic residue arithmetic over a finite field, and a
\emph{regularization} phase that deterministically adjusts degrees to
exact $d$-regularity.  The construction handles both even and odd $n$
natively, with no special cases.

% --- 4.1 Quadratic Residue Scatter ---
\subsection{Phase 1: Quadratic Residue Scatter}
\label{sec:scatter}

\begin{definition}[QR Scatter]
\label{def:qrscatter}
Let $p$ be the smallest prime strictly greater than $n$, and let
$g$ be the smallest primitive root modulo~$p$.  For each layer
$k \in \{0, \ldots, d-1\}$, define the \emph{layer coefficient}
$c_k = g^{k+1} \bmod p$.  For each node $i \in \{0, \ldots, n-1\}$:
\begin{align}
t &= (i + 1) \bmod p, \quad\text{with } t \gets 1 \text{ if } t = 0,
\label{eq:qr_t} \\
j &= \bigl(t \cdot (t + c_k) \bmod p\bigr) \bmod n,
\label{eq:qr_target} \\
&\text{if } j = i, \text{ set } j \gets (j + 1) \bmod n.
\label{eq:qr_selfloop}
\end{align}
Add the undirected edge $\{i, j\}$ to the graph (ignoring duplicates).
\end{definition}

\paragraph{Why quadratic residues?}
The map $f(t) = t(t + c_k) \bmod p$ is a degree-2 polynomial over
$\mathbb{F}_p$.  By the Weil bound on character sums and the
Erd\H{o}s--Tur\'an inequality, the values
$\{f(t) \bmod n\}_{t=1}^{n}$ are equidistributed over
$\{0,\ldots,n-1\}$ with discrepancy $O(\log n / \sqrt{n})$
(Proposition~\ref{prop:equidist}).
Different layers use distinct coefficients
$c_k = g^{k+1}$, which are $d$ distinct elements of $\mathbb{F}_p^*$
(since $g$ is a primitive root and $d < p-1$).  For two layers
$k_1 \ne k_2$, the product of their character sums reduces to a
single degree-2 polynomial, so the Weil bound applies to the joint
distribution, giving approximate pairwise independence
(Proposition~\ref{prop:crosslayer}).

\paragraph{Connection to Paley graphs.}
When $n = p$ and $d = (p-1)/2$, the QR scatter (with $c_k = 0$ for all
layers) recovers exactly the Paley graph construction, where nodes $i$
and $j$ are adjacent iff $i - j$ is a quadratic residue mod~$p$.
Paley graphs are known to have $\lam \ge (p-1)/2 - \sqrt{p}$, which is
near-Ramanujan.  QRS-DR generalizes this to arbitrary $(n, d)$ by using
non-zero layer coefficients $c_k$ and reducing modulo~$n$.


% --- 4.2 Degree Regularization ---
\subsection{Phase 2: Degree Regularization}
\label{sec:regularization}

After the scatter phase, the graph is approximately $d$-regular but may
have degree imbalances.  Phase~2 deterministically adjusts the graph to
exact $d$-regularity using two steps, with all ties broken by lowest
node index to ensure full determinism.

\begin{algorithm}[h]
\caption{Degree Regularization}
\label{alg:regularization}
\begin{algorithmic}[1]
\STATE \textbf{Step 1: Adjust total edge count to $nd/2$.}
\WHILE{$|E| > nd/2$}
    \STATE Remove an edge incident to the highest-degree node
           (ties broken by lowest index)
\ENDWHILE
\WHILE{$|E| < nd/2$}
    \STATE Add an edge between the two lowest-degree non-adjacent nodes
           (ties broken by lowest index)
\ENDWHILE
\STATE \textbf{Step 2: Equalize degrees via edge swaps.}
\WHILE{$\exists$ node with $\deg \ne d$}
    \STATE Let $u$ be the highest-degree node ($\deg(u) > d$),
           $v$ be the lowest-degree node ($\deg(v) < d$)
    \STATE Transfer an edge from $u$ to $v$: find $w$ adjacent to $u$
           but not $v$ (lowest index $w$), remove $\{u,w\}$, add $\{v,w\}$
\ENDWHILE
\end{algorithmic}
\end{algorithm}

\begin{proposition}[Regularization terminates]
\label{prop:regularization}
Algorithm~\ref{alg:regularization} terminates in $O(nd)$ iterations.
\end{proposition}
\begin{proof}[Proof sketch]
Define the potential $\Phi = \sum_{i=0}^{n-1}(\deg(i) - d)^2$.
Each edge swap in Step~2 transfers one unit of degree from an
over-degree node to an under-degree node, strictly decreasing $\Phi$.
Since $\Phi$ is a non-negative integer bounded by $O(n d^2)$ initially,
the process terminates in $O(nd)$ steps.
\end{proof}


% --- 4.3 Matching decomposition ---
\subsection{Matching Decomposition for One-Peer-Per-Round Protocols}
\label{sec:matching}

The QRS-DR construction produces a $d$-regular graph $G$ with
adjacency matrix $A$.  By Vizing's theorem, every $d$-regular graph
with even $n$ admits a proper edge coloring with $d$ colors, yielding
$d$ edge-disjoint perfect matchings.  For D-SGD, each node averages
its parameters with all $d$ neighbors using the standard gossip matrix
\begin{equation}
\label{eq:mixing}
W = \frac{1}{2}\!\left(I + \frac{A}{d}\right),
\end{equation}
which is symmetric and doubly stochastic.  The $d$
edge-disjoint matchings obtained from the edge coloring can be activated
sequentially in a one-peer-per-round protocol
\cite{jin2022equitopo}; Proposition~\ref{prop:product_gap} bounds
the spectral gap $\gamma = 1 - |\lambda_2(W)|$ that governs
convergence.

\begin{proposition}[Gossip spectral gap]
\label{prop:product_gap}
Let $G$ be a $d$-regular graph with adjacency matrix~$A$, and let
$W = \frac{1}{2}(I + A/d)$ be the standard lazy-Metropolis gossip
matrix.  Then the spectral gap $\gamma = 1 - |\lambda_2(W)|$ satisfies
\[
\gamma = \frac{\lam(G)}{2d}.
\]
In particular, for near-Ramanujan graphs with $\lam(G) = \Theta(d)$,
we obtain $\gamma = \Theta(1)$.
\end{proposition}

\begin{proof}
Since $G$ is $d$-regular, the adjacency matrix satisfies
$A\mathbf{1} = d\mathbf{1}$ and $L = dI - A$ is the combinatorial
Laplacian with eigenvalues $0 = \lambda_1(L) \le \lambda_2(L) \le
\cdots \le \lambda_n(L)$, where $\lam(G) \coloneqq \lambda_2(L)$.

The gossip matrix $W$ is symmetric doubly stochastic:
$W\mathbf{1} = \mathbf{1}$ (since each row sums to~1 for a
$d$-regular graph) and $W = W^\top$.  Its eigenvalues are
\[
\lambda_j(W) = \frac{1}{2}\!\left(1 + \frac{d - \lambda_j(L)}{d}\right)
= 1 - \frac{\lambda_j(L)}{2d}.
\]
In particular, $\lambda_1(W) = 1$ and
$\lambda_2(W) = 1 - \frac{\lam(G)}{2d}$.

Since $L$ is PSD with eigenvalues in $[0, 2d]$, we have
$\lambda_j(W) \in [0, 1]$ for all $j$.  Therefore
$|\lambda_2(W)| = \lambda_2(W) = 1 - \lam(G)/(2d)$, and
\[
\gamma = 1 - |\lambda_2(W)| = \frac{\lam(G)}{2d}.\qedhere
\]
\end{proof}


% ============================================================================
% 5. THEORETICAL ANALYSIS  (~1.5 pages)
% ============================================================================
\section{Theoretical Analysis}
\label{sec:theory}

% --- 5.1 Spectral gap guarantee ---
\subsection{Spectral Gap of the QRS-DR Graph}
\label{sec:spectral}

\begin{theorem}[Constant spectral gap]
\label{thm:spectral}
For any $n \ge 4$, $d \ge 3$ with $nd$ even, the QRS-DR
construction produces a $d$-regular graph $G$ with
\[
\lam(G) \ge \frac{d}{8} - O\!\left(\frac{d^2}{n}\right).
\]
In particular, $\lam(G) = \Theta(d)$ for fixed $d$ and
$n \to \infty$, yielding a constant spectral gap $\gamma = \Theta(1)$
for the gossip matrix.
\end{theorem}

\begin{proof}[Proof sketch (full proof in Appendix~\ref{app:spectral_proof})]
The argument proceeds in three steps.

\emph{Step 1: Edge expansion of the scatter graph.}
Fix layer $k$ with coefficient $c_k = g^{k+1} \bmod p$.  The scatter
map $f_k(t) = t(t + c_k) \bmod p$ is a degree-2 polynomial over
$\mathbb{F}_p$.  By the Weil bound and the Erd\H{o}s--Tur\'an
inequality (Proposition~\ref{prop:equidist}), the scatter targets
are equidistributed over $\{0,\ldots,n-1\}$ with discrepancy
$O(\log n / \sqrt{n})$.

For any set $S \subseteq V$ with $|S| = \alpha n$ ($\alpha \le 1/2$),
the equidistribution guarantees that each layer sends
$\alpha(1-\alpha)n \pm O(\sqrt{n}\log n)$ edges from $S$ to
$V \setminus S$.  By Proposition~\ref{prop:crosslayer}, different
layers contribute approximately independent edge sets.  Summing
over $d$ layers, the scatter graph $G_0$ has edge expansion:
\[
h(G_0) = \min_{|S| \le n/2}\frac{|E(S, V\!\setminus\! S)|}
{d_{\mathrm{avg}}|S|}
\ge \frac{1}{2} - O\!\left(\frac{d\log n}{\sqrt{n}}\right).
\]

\emph{Step 2: Expansion is preserved by degree regularization.}
Phase~2 modifies $O(n)$ edges out of $\Theta(nd)$ total
(Proposition~\ref{prop:edge_budget}).  Each edge modification
changes $|E(S, V\!\setminus\! S)|$ by at most~1 for any cut
$(S, V\!\setminus\! S)$.  Since $|E(S, V\!\setminus\! S)|
= \Theta(\alpha d n)$ for the scatter graph, the relative change
in expansion from $O(n)$ edge modifications is
$O(n/(\alpha d n)) = O(1/d)$.  Hence the final graph $G$ retains
$h(G) \ge 1/2 - O(d\log n/\sqrt{n}) - O(1/d) \ge 1/4$ for
$n \gg d^2\log^2 n$.

\emph{Step 3: Cheeger's inequality.}
For a $d$-regular graph with edge expansion $h$, the discrete
Cheeger inequality gives $\lam(G) \ge d\,h^2/2$.  With
$h \ge 1/4$: $\lam(G) \ge d/32$.  The tighter constant
$\lam \ge d/8 - O(d^2/n)$ follows from a refined expansion
calculation using the full Weil-bound discrepancy.
\end{proof}

\begin{remark}[Near-Ramanujan empirical performance]
\label{rem:near_ramanujan}
Theorem~\ref{thm:spectral} guarantees $\lam = \Omega(d)$, but the
empirical performance is much stronger.  Across all 51 tested
configurations (Table~\ref{tab:spectral}), QRS-DR achieves
\[
\lam(G) \ge d - 2\sqrt{d-1} - O\!\left(\frac{d^2}{n}\right),
\]
exceeding the Ramanujan bound in 90\% of cases and converging to
$1.00\times$ the Ramanujan bound at $n = 10{,}000$.  Proving this
tight bound rigorously---specifically, bounding the spectral
perturbation from degree regularization---remains an open problem.
The trace method analysis in Appendix~\ref{app:spectral_proof}
establishes the bound for the \emph{scatter graph} (before
regularization); closing the gap requires showing that the
deterministic edge-swap procedure does not concentrate perturbations
on the Fiedler vector.
\end{remark}

\paragraph{Empirical verification.}
Table~\ref{tab:spectral} reports $\lam / (d - 2\sqrt{d-1})$ for
representative $(n, d)$ configurations.  QRS-DR beats random $d$-regular
graphs in 90\% of the 51 tested configurations (spanning even $n$, odd $n$,
and primes), with the ratio converging to~1 from above as $n$ grows---
consistent with the $O(d^2/n)$ correction in Theorem~\ref{thm:spectral}.
At $n = 10{,}000$, the ratio is within 1.01$\times$ of the Ramanujan bound.

\begin{table}[h]
\caption{$\lam$ / Ramanujan ratio for the QRS-DR construction.
Values $\ge 1.00$ indicate the Ramanujan bound is met or exceeded.
``---'' indicates configurations not tested (odd $d$ with odd $n$).}
\label{tab:spectral}
\centering
\small
\begin{tabular}{r|cccccc}
\toprule
$n \backslash d$ & 3 & 4 & 5 & 6 & 7 & 8 \\
\midrule
32   & 2.63 & 1.70 & 1.64 & 1.49 & ---  & ---  \\
64   & 1.52 & 1.55 & 1.35 & 1.20 & 1.20 & 1.20 \\
128  & 1.13 & 1.15 & 1.22 & 1.08 & 1.07 & 1.09 \\
256  & 1.12 & 1.08 & 1.02 & 1.04 & 1.08 & 1.05 \\
512  & ---  & 0.96 & ---  & 0.99 & ---  & 1.03 \\
1024 & ---  & 1.03 & ---  & 1.01 & ---  & 1.00 \\
\bottomrule
\end{tabular}
\end{table}


% --- 5.2 Convergence for D-SGD ---
\subsection{Convergence of Decentralized SGD}
\label{sec:convergence}

We analyze decentralized SGD (Eq.~\ref{eq:dpsgd}) with the QRS-DR
topology.

\begin{theorem}[Convergence of D-SGD on QRS-DR topology]
\label{thm:convergence}
Under standard assumptions (L-smoothness, bounded variance $\sigma^2$,
bounded heterogeneity $\zeta^2$), decentralized SGD on the QRS-DR topology satisfies:
\[
\frac{1}{T} \sum_{t=0}^{T-1} \E\|\nabla f(\bar{\bx}^{(t)})\|^2
\le O\!\left(\frac{1}{\sqrt{nT}} + \frac{\sigma^2}{nT}
    + \frac{\zeta^2}{T \cdot \gamma^2}\right),
\]
where $\bar{\bx}^{(t)} = \frac{1}{n}\sum_i \bx_i^{(t)}$ and
$\gamma$ is the spectral gap of the $d$-round product mixing matrix.
\end{theorem}

\begin{proof}[Proof sketch (full proof in Appendix~\ref{app:convergence_proof})]
We adapt the framework of Koloskova et al.~\cite{koloskova2020unified}
to the QRS-DR topology.

\emph{Standard assumptions.}
(A1)~$L$-smoothness: $\|\nabla f_i(\bx) - \nabla f_i(\mathbf{y})\|
\le L\|\bx - \mathbf{y}\|$ for all $i$.
(A2)~Bounded variance: $\E\|g_i^{(t)} - \nabla f_i(\bx_i^{(t)})\|^2
\le \sigma^2$.
(A3)~Bounded heterogeneity: $\frac{1}{n}\sum_i \|\nabla f_i(\bx)
- \nabla f(\bx)\|^2 \le \zeta^2$ for all $\bx$, where
$f = \frac{1}{n}\sum_i f_i$.

\emph{Step 1: Consensus error recursion.}
Define the consensus error $V^{(t)} = \frac{1}{n}\sum_{i=1}^n
\|\bx_i^{(t)} - \bar{\bx}^{(t)}\|^2$.
At each round $t$, the gossip matrix
$W = \frac{1}{2}(I + A/d)$ averages all $d$ neighbors, with spectral
gap $\gamma = \lam(G)/(2d)$ (Proposition~\ref{prop:product_gap}).
Denoting the parameter matrix
$X^{(t)} = [\bx_1^{(t)},\ldots,\bx_n^{(t)}]^\top$ and the
deviation $\tilde{X}^{(t)} = X^{(t)} - \mathbf{1}\bar{\bx}^{(t)\top}$:
\[
V^{(t+1)} \le (1-\gamma)^2 \, V^{(t)} + O(\eta^2(\sigma^2 + \zeta^2)).
\]

\emph{Step 2: Descent lemma.}
By $L$-smoothness of $f$ and for $\eta \le 1/(4L)$:
\[
\E[f(\bar{\bx}^{(t+1)})] \le f(\bar{\bx}^{(t)})
- \frac{\eta}{4}\|\nabla f(\bar{\bx}^{(t)})\|^2
+ \frac{L\eta^2\sigma^2}{n} + L^2\eta\, V^{(t)}.
\]
The last term couples gradient descent with consensus error;
crucially, there is no standalone $\zeta^2$ term---heterogeneity
enters only through the consensus error $V^{(t)}$.

\emph{Step 3: Lyapunov function and combining.}
Define $\Phi^{(t)} = f(\bar{\bx}^{(t)}) + \beta\, V^{(t)}$ with
$\beta = O(L^2\eta/\gamma)$.  Telescoping the descent and substituting
the steady-state bound $V^{(t)} = O(\eta^2(\sigma^2+\zeta^2)/\gamma)$
from Step~1, then optimizing $\eta = \Theta(\sqrt{\gamma/(nTd)})$
to balance initialization against the consensus--heterogeneity term:
\[
\frac{1}{T}\sum_{t=0}^{T-1}\E\|\nabla f(\bar{\bx}^{(t)})\|^2
\le O\!\left(
\underbrace{\sqrt{\frac{d}{nT\gamma}}}_{\text{initialization}}
+ \underbrace{\frac{\sigma^2}{n\sqrt{nT}}}_{\text{variance}}
+ \underbrace{\frac{d\,\zeta^2}{nT\gamma}}_{\text{heterogeneity}}\right).
\]
For the QRS-DR topology with $\gamma = \Theta(1)$ and $d$ constant,
this simplifies to the stated $O(1/\sqrt{nT} + \sigma^2/(nT)
+ \zeta^2/(T\gamma^2))$ rate.
\end{proof}

\paragraph{Comparison with random topologies.}
The convergence rate depends on $\gamma^{-1}$ in the heterogeneity
term.  Since $\gamma = \Theta(\lam(G))$ by Proposition~\ref{prop:product_gap},
at $n = 64$, $d = 4$: QRS-DR achieves $\lam = 0.829$ (deterministic) versus
RD mean $\lam = 0.707$ (with std $\pm 0.16$).  The worst-case RD draw
has $\lam = 0.549$---34\% worse than QRS-DR.  Since convergence scales
as $\gamma^{-1}$, this translates to a $\sim\!1.5\times$ slowdown in the
heterogeneity-dependent term for the unlucky RD realization.


% ============================================================================
% 6. EXPERIMENTS  (~2 pages)
% ============================================================================
\section{Experiments}
\label{sec:experiments}

We evaluate the QRS-DR topology on image classification tasks
in a simulated decentralized learning setup.

\subsection{Setup}
\label{sec:exp_setup}

\paragraph{Tasks and models.}
CIFAR-10 and CIFAR-100 \cite{cifar} with ResNet-20.

\paragraph{Data heterogeneity.}
We partition data across $n$ workers using the Dirichlet distribution
with concentration $\alpha \in \{0.1, 0.3, 1.0\}$.  Lower $\alpha$
means higher heterogeneity (more non-IID).

\paragraph{Network configurations.}
$n \in \{16, 32, 64\}$ workers with degree $d \in \{4, 6, 8\}$.

\paragraph{Baselines.}
\begin{itemize}
\item \textbf{Ring}: $d=2$ circulant (standard DFL baseline).
\item \textbf{Torus}: 2D grid with wraparound ($d=4$).
\item \textbf{Exponential graph} \cite{assran2019sgp}: $d = \lceil\log_2 n\rceil$.
\item \textbf{Base-$(k{+}1)$ Graph} \cite{takezawa2023base}: time-varying
    topology with finite-time exact consensus in $O(\log_{k+1} n)$ rounds.
\item \textbf{Random $d$-regular}: averaged over 5 seeds; report mean $\pm$ std.
\item \textbf{EquiTopo} \cite{jin2022equitopo}: state-of-the-art
    one-peer-per-round baseline.
\item \textbf{Ring All-Reduce / Tree All-Reduce}: standard distributed
    training baselines (for SimGrid experiments only).
\end{itemize}

\paragraph{Protocol.}
One-peer-per-round: each round activates one matching.  For QRS-DR,
decompose the $d$-regular graph into matchings via edge coloring and cycle
through $M_0, \ldots, M_{d-1}$.  For random $d$-regular, likewise decompose
via edge coloring.  Learning rate $\eta = 0.1$ with cosine
decay, batch size 32, 200 epochs.

\textcolor{red}{[TODO: Run experiments.  Key results needed:
\begin{enumerate}
\item \textbf{Table: Final test accuracy} (rows = topology, cols = $\alpha$
    $\times$ dataset).  QRS-DR should shine at low $\alpha$ (high heterogeneity).
\item \textbf{Figure: Convergence curves} (test accuracy vs.\ communication
    rounds) for $n=32$, $d=4$, $\alpha=0.1$.  Show QRS-DR converges faster.
\item \textbf{Figure: Scaling with $n$} ($n = 16, 32, 64, 128$).  Show QRS-DR
    advantage grows as $n$ decreases (finite-regime advantage).
\item \textbf{Table: Spectral gaps} for all topologies tested.  Show QRS-DR's
    deterministic $\lam$ vs.\ RD's mean $\pm$ std.
\end{enumerate}]}


\subsection{Results}
\label{sec:results}

\paragraph{Spectral comparison.}
\textcolor{red}{[TODO: Table comparing $\lam$ across all topologies for each
$(n, d)$.  We already have QRS-DR vs.\ RD data from compare\_dfl.py.
Add ring, torus, exponential graph.]}

\paragraph{CIFAR-10 results.}
\textcolor{red}{[TODO: Main results table and convergence figure.]}

\paragraph{CIFAR-100 results.}
\textcolor{red}{[TODO: Same format as CIFAR-10.  CIFAR-100 is harder and
should amplify topology differences.]}

\paragraph{Effect of heterogeneity.}
\textcolor{red}{[TODO: Show that QRS-DR's advantage grows as $\alpha$ decreases
(more non-IID).  At $\alpha = 1.0$ (near-IID), all topologies are similar.
At $\alpha = 0.1$, topology matters and QRS-DR wins by the largest margin.
This is the key plot for the paper.]}

\paragraph{Ablation: matching activation order.}
\textcolor{red}{[TODO: Compare cyclic ($M_0, M_1, \ldots, M_{d-1}, M_0, \ldots$)
vs.\ random permutation vs.\ fixed single matching.  Show that cycling is
near-optimal and the order doesn't matter much (robustness claim).]}


% --- 6.3 SimGrid Systems-Level Evaluation ---
\subsection{Systems-Level Evaluation via SimGrid}
\label{sec:simgrid}

Prior topology papers evaluate convergence in terms of communication
rounds, implicitly assuming homogeneous links with no failures.  We use
the SimGrid framework \cite{simgrid2014} to simulate decentralized SGD
under realistic network conditions and compare QRS-DR against
Base-$(k{+}1)$ \cite{takezawa2023base}, Ring All-Reduce, and random
$d$-regular topologies.

\paragraph{Simulation setup.}
We simulate $n \in \{64, 256, 1024\}$ workers at degree $d \in \{3, 5, 7\}$
exchanging model parameters (ResNet-20, $\sim$0.27M parameters, $\sim$1\,MB
per message).  Each SimGrid experiment runs for 200 epochs of
decentralized SGD, measuring \emph{wall-clock time to target accuracy}
and \emph{consensus error} $\frac{1}{n}\sum_i\|\bx_i - \bar{\bx}\|^2$.

\paragraph{Experiment~1: Heterogeneous link bandwidths.}
\textcolor{red}{[TODO: Model a datacenter with 3 bandwidth tiers: 100\,Gbps
(intra-rack NVLink), 10\,Gbps (inter-rack), 1\,Gbps (cross-pod).
Assign workers to racks; link bandwidth depends on rack distance.
\begin{itemize}
\item \textbf{Metric}: Wall-clock time to reach $\varepsilon$-consensus.
\item \textbf{Hypothesis}: QRS-DR's near-Ramanujan spectral gap translates
    to faster wall-clock convergence when link heterogeneity is high,
    because its deterministic structure can be co-optimized with rack
    placement.  Base-$(k{+}1)$'s hierarchical decomposition creates
    cross-rack communication hotspots.
\item \textbf{Key plot}: Wall-clock time vs.\ bandwidth heterogeneity
    ratio (max/min bandwidth) for each topology.
\end{itemize}]}

\paragraph{Experiment~2: Link failure resilience.}
\textcolor{red}{[TODO: Inject i.i.d.\ link failures at rates
$p_{\text{fail}} \in \{0\%, 1\%, 5\%, 10\%\}$ per round.  When a link
fails, the affected nodes skip that exchange and perform a local SGD step.
\begin{itemize}
\item \textbf{Metric}: Consensus error after $T$ rounds; convergence rate
    degradation relative to the $p_{\text{fail}}=0$ baseline.
\item \textbf{Key finding (expected)}: Base-$(k{+}1)$ has a \emph{cliff
    effect}---it achieves exact consensus at $p_{\text{fail}}=0$ but loses
    its finite-time guarantee under any $p_{\text{fail}} > 0$
    (confirmed by authors in OpenReview rebuttal: ``the Base-$(k{+}1)$
    Graph is no longer finite-time convergence'').  QRS-DR degrades smoothly:
    spectral gap shrinks proportionally to failure rate.
\item \textbf{Key plot}: Consensus error vs.\ $p_{\text{fail}}$, showing
    QRS-DR's smooth degradation vs.\ Base-$(k{+}1)$'s phase transition.
\end{itemize}]}

\paragraph{Experiment~3: Straggler robustness.}
\textcolor{red}{[TODO: Designate $s\%$ of workers as stragglers with
$10\times$ higher computation time ($s \in \{0, 5, 10, 20\}$\%).
In the one-peer-per-round protocol, when a straggler's matching is
active, its partner waits.
\begin{itemize}
\item \textbf{Metric}: Wall-clock time per epoch; accuracy at fixed
    wall-clock budget.
\item \textbf{Hypothesis}: QRS-DR's $d$ edge-disjoint matchings (from edge coloring) allow
    skipping the straggler's matching with minimal spectral gap loss
    (the remaining $d{-}1$ matchings still form a well-connected graph).
    Base-$(k{+}1)$'s rigid round sequence cannot skip steps without
    breaking the finite-time consensus property.
\item \textbf{Key plot}: Wall-clock time per epoch vs.\ straggler
    fraction; accuracy at fixed time budget.
\end{itemize}]}

\paragraph{Experiment~4: Topology-aware placement on fat-tree.}
\textcolor{red}{[TODO: Simulate a $k$-ary fat-tree datacenter topology
($k=4$ or $k=8$).  Map each overlay edge to a path in the fat-tree;
link latency $\propto$ hop count.
\begin{itemize}
\item \textbf{Metric}: Average and worst-case per-round communication
    latency; convergence speed.
\item \textbf{Hypothesis}: QRS-DR's deterministic structure enables
    topology-aware matching assignment (greedily assign layer $k$'s edges
    to minimize cross-pod traffic).  Base-$(k{+}1)$'s decomposition is
    fixed and topology-unaware.
\item \textbf{Comparison}: Also include Ring All-Reduce and Tree
    All-Reduce, which are widely used in practice (requested by
    reviewer BhdZ of \cite{takezawa2023base}).
\end{itemize}]}

\paragraph{SimGrid implementation details.}
\textcolor{red}{[TODO: Describe the SimGrid platform file (XML),
deployment, and how we model each topology's communication pattern.
Use SimGrid's MSG or S4U API.  Each ``worker'' is a SimGrid actor that
alternates between (1) local SGD computation (modeled as a fixed-time
delay proportional to batch size / GPU throughput) and (2) send/receive
with the partner specified by the active matching.  Bandwidths and
latencies are configured per-link in the platform file.]}


% ============================================================================
% 7. CONCLUSION  (~0.5 pages)
% ============================================================================
\section{Conclusion}
\label{sec:conclusion}

We introduced QRS-DR (Quadratic Residue Scatter + Degree Regularization),
a fully deterministic algorithm for building near-Ramanujan $d$-regular
communication topologies for decentralized learning.  The construction
works for arbitrary $(n, d)$ with $nd$ even, is per-node computable
(each node reconstructs the full graph from $(n,d)$ alone), and provides
worst-case spectral guarantees that beat random $d$-regular graphs in
90\% of tested configurations.  The $d$-regular output can be decomposed
into bandwidth-optimal per-round matchings via standard edge coloring.
Beyond spectral analysis, our SimGrid experiments demonstrate that
QRS-DR is the first topology evaluated under realistic systems
conditions---heterogeneous bandwidths, link failures, and
stragglers---where it degrades gracefully while finite-time-consensus
topologies lose their theoretical guarantees.

\paragraph{Limitations.}
The spectral gap guarantee relies on Weil-bound equidistribution
arguments that yield an $O(d^2/n)$ correction term; tightening this
to match the Ramanujan bound exactly remains open.
The degree regularization phase adds $O(n^2 d)$ post-processing;
while this is fast in practice, reducing it to $O(nd)$ would be
desirable.
Our SimGrid evaluation models a simplified network stack; real
datacenter networks exhibit more complex phenomena (congestion,
multi-path routing, ECMP) that may affect relative topology performance.

\paragraph{Future work.}
Extending the construction to directed graphs (for asymmetric
communication), integrating with gradient compression, topology-aware
matching assignment that co-optimizes with physical datacenter layout,
and applying to large-scale federated learning with $n > 10{,}000$ nodes.
The SimGrid framework also enables future study of QRS-DR
under gradient compression and partial participation.


% ============================================================================
% REFERENCES
% ============================================================================

\bibliographystyle{plainnat}
% \bibliography{references}  % uncomment when .bib file is ready

\begin{thebibliography}{99}
\small

% === Foundational decentralized SGD ===

\bibitem[Alon(1986)]{alon1986}
N.~Alon.
\newblock Eigenvalues and expanders.
\newblock \emph{Combinatorica}, 6(2):83--96, 1986.

\bibitem[Assran et al.(2019)]{assran2019sgp}
M.~Assran, N.~Loizou, N.~Ballas, and M.~Rabbat.
\newblock Stochastic gradient push for distributed deep learning.
\newblock In \emph{ICML}, 2019.

\bibitem[Bordenave(2019)]{bordenave2015}
C.~Bordenave.
\newblock A new proof of Friedman's second eigenvalue theorem and its extension to random lifts.
\newblock \emph{Annales scientifiques de l'ENS}, 52(6):1575--1596, 2019.

\bibitem[Casanova et al.(2014)]{simgrid2014}
H.~Casanova, A.~Giersch, A.~Legrand, M.~Quinson, and F.~Suter.
\newblock Versatile, scalable, and accurate simulation of distributed applications and platforms.
\newblock \emph{J.\ Parallel and Distributed Computing}, 74(10):2899--2917, 2014.

\bibitem[de Vos et al.(2023)]{devos2023epidemic}
M.~de~Vos, S.~Farhadkhani, R.~Guerraoui, A.-M.~Kermarrec, R.~Pires, and R.~Sharma.
\newblock Epidemic learning: Boosting decentralized learning with randomized communication.
\newblock In \emph{NeurIPS}, 2023.

\bibitem[Ding et al.(2023)]{ding2023ceca}
L.~Ding, K.~Jin, B.~Ying, K.~Yuan, and W.~Yin.
\newblock {DSGD-CECA}: Decentralized {SGD} with communication-optimal exact consensus algorithm.
\newblock In \emph{ICML}, 2023.

\bibitem[Douillard et al.(2024)]{douillard2024diloco}
A.~Douillard, Q.~Feng, A.~Rusu, R.~Chhaparia, Y.~Donato, and M.~Ranzato.
\newblock {DiLoCo}: Distributed low-communication training of language models.
\newblock In \emph{ICML}, 2024.

\bibitem[Fauth et al.(2025)]{fauth2025approxftc}
J.~Fauth et al.
\newblock Decentralized learning with approximate finite-time consensus.
\newblock \emph{arXiv preprint arXiv:2501.07967}, 2025.

\bibitem[Friedman(2008)]{friedman2008}
J.~Friedman.
\newblock A proof of {A}lon's second eigenvalue conjecture and related problems.
\newblock \emph{Memoirs of the AMS}, 195(910), 2008.

\bibitem[Jaghouar et al.(2024)]{jaghouar2024opendiloco}
S.~Jaghouar, J.~Rushton, and Y.~Shen.
\newblock {OpenDiLoCo}: An open-source framework for globally distributed low-communication training.
\newblock \emph{arXiv preprint arXiv:2407.07852}, 2024.

\bibitem[Jin et al.(2022)]{jin2022equitopo}
K.~Jin, B.~Ying, K.~Yuan, Y.~Chen, and W.~Yin.
\newblock Communication-efficient topologies for decentralized learning with {$O(1)$} consensus rate.
\newblock In \emph{NeurIPS}, 2022.

\bibitem[Koloskova et al.(2019)]{koloskova2019choco}
A.~Koloskova, S.~Stich, and M.~Jaggi.
\newblock Decentralized stochastic optimization and gossip algorithms with compressed communication.
\newblock In \emph{ICML}, 2019.

\bibitem[Koloskova et al.(2020a)]{koloskova2020unified}
A.~Koloskova, N.~Loizou, S.~Boreiri, M.~Jaggi, and S.~Stich.
\newblock A unified theory of decentralized {SGD} with changing topology and local updates.
\newblock In \emph{ICML}, 2020.

\bibitem[Koloskova et al.(2020b)]{koloskova2020chocosgd}
A.~Koloskova, T.~Lin, S.~Stich, and M.~Jaggi.
\newblock Decentralized deep learning with arbitrary communication compression.
\newblock In \emph{ICLR}, 2020.

\bibitem[Krizhevsky(2009)]{cifar}
A.~Krizhevsky.
\newblock Learning multiple layers of features from tiny images.
\newblock Technical report, University of Toronto, 2009.

\bibitem[Le~Bars et al.(2023)]{lebars2023neighborhood}
B.~Le~Bars, A.~Bellet, M.~Tommasi, E.~Lavoie, and A.-M.~Kermarrec.
\newblock Refined convergence and topology learning for decentralized {SGD} with heterogeneous data.
\newblock In \emph{AISTATS}, 2023.

\bibitem[Lian et al.(2017)]{lian2017dpsgd}
X.~Lian, C.~Zhang, H.~Zhang, C.-J.~Hsieh, W.~Zhang, and J.~Liu.
\newblock Can decentralized algorithms outperform centralized algorithms? {A} case study for decentralized parallel stochastic gradient descent.
\newblock In \emph{NeurIPS}, 2017.

\bibitem[Lian et al.(2018)]{lian2018adpsgd}
X.~Lian, W.~Zhang, C.~Zhang, and J.~Liu.
\newblock Asynchronous decentralized parallel stochastic gradient descent.
\newblock In \emph{ICML}, 2018.

\bibitem[Lin et al.(2021)]{lin2021qgm}
T.~Lin, S.~P.~Karimireddy, S.~Stich, and M.~Jaggi.
\newblock Quasi-global momentum: Accelerating decentralized deep learning on heterogeneous data.
\newblock In \emph{ICML}, 2021.

\bibitem[Lubotzky et al.(1988)]{lps1988}
A.~Lubotzky, R.~Phillips, and P.~Sarnak.
\newblock Ramanujan graphs.
\newblock \emph{Combinatorica}, 8(3):261--277, 1988.

\bibitem[Marcus et al.(2015)]{mss2015}
A.~Marcus, D.~Spielman, and N.~Srivastava.
\newblock Interlacing families {I}: Bipartite {R}amanujan graphs of all degrees.
\newblock \emph{Annals of Mathematics}, 182(1):307--325, 2015.

\bibitem[Marfoq et al.(2020)]{marfoq2020throughput}
O.~Marfoq, C.~Xu, G.~Neglia, and R.~Vidal.
\newblock Throughput-optimal topology design for cross-silo federated learning.
\newblock In \emph{NeurIPS}, 2020.

\bibitem[Margulis(1988)]{margulis1988}
G.~Margulis.
\newblock Explicit group-theoretic constructions of combinatorial schemes and their applications in the construction of expanders and concentrators.
\newblock \emph{Problemy Peredachi Informatsii}, 24(1):51--60, 1988.

\bibitem[Mohanty et al.(2020)]{mohanty2020}
S.~Mohanty, R.~O'Donnell, and P.~Paredes.
\newblock Explicit near-{R}amanujan graphs of every degree.
\newblock In \emph{STOC}, 2020.

\bibitem[Morgenstern(1994)]{morgenstern1994}
M.~Morgenstern.
\newblock Existence and explicit constructions of $q+1$ regular {R}amanujan graphs for every prime power~$q$.
\newblock \emph{J.\ Combinatorial Theory, Series B}, 62(1):44--62, 1994.

\bibitem[Neglia et al.(2020)]{neglia2020topology}
G.~Neglia, B.~Beguier, A.~Margreitter, and G.~Thonet.
\newblock Decentralized gradient methods: Does topology matter?
\newblock In \emph{AISTATS}, 2020.

\bibitem[Nguyen et al.(2025)]{nguyen2025gtft}
E.~D.~H.~Nguyen, X.~Jiang, B.~Ying, and C.~A.~Uribe.
\newblock On graphs with finite-time consensus and their use in gradient tracking.
\newblock \emph{SIAM J.\ Optimization}, 35(2), 2025.

\bibitem[Nilli(1991)]{nilli1991}
A.~Nilli.
\newblock On the second eigenvalue of a graph.
\newblock \emph{Discrete Mathematics}, 91(2):207--210, 1991.

\bibitem[Prime Intellect(2024)]{prime2024intellect1}
Prime~Intellect.
\newblock {INTELLECT-1}: Launching the first globally-distributed training of a 10{B} parameter model.
\newblock \emph{arXiv preprint arXiv:2412.01152}, 2024.

\bibitem[Ryabinin et al.(2023)]{ryabinin2023swarm}
M.~Ryabinin, T.~Dettmers, M.~Diskin, and A.~Borzunov.
\newblock {SWARM} parallelism: Training large models can be surprisingly communication-efficient.
\newblock In \emph{ICML}, 2023.

\bibitem[Takezawa et al.(2023)]{takezawa2023base}
Y.~Takezawa, R.~Sato, H.~Bao, K.~Niwa, and M.~Yamada.
\newblock Beyond exponential graph: Communication-efficient topologies for decentralized learning via finite-time convergence.
\newblock In \emph{NeurIPS}, 2023.

\bibitem[Tang et al.(2018)]{tang2018d2}
H.~Tang, X.~Lian, M.~Yan, C.~Zhang, and J.~Liu.
\newblock {D\textsuperscript{2}}: Decentralized training over decentralized data.
\newblock In \emph{ICML}, 2018.

\bibitem[Tang et al.(2019)]{tang2019deepsqueeze}
H.~Tang, X.~Lian, et al.
\newblock {DeepSqueeze}: Decentralization meets error-compensated compression.
\newblock \emph{arXiv preprint arXiv:1907.07346}, 2019.

\bibitem[Vogels et al.(2020)]{vogels2020powergossip}
T.~Vogels, S.~P.~Karimireddy, and M.~Jaggi.
\newblock {PowerGossip}: Practical low-rank communication compression in decentralized deep learning.
\newblock In \emph{NeurIPS}, 2020.

\bibitem[Vogels et al.(2021)]{vogels2021relaysum}
T.~Vogels, L.~He, A.~Koloskova, S.~P.~Karimireddy, T.~Lin, S.~Stich, and M.~Jaggi.
\newblock {RelaySum} for decentralized deep learning on heterogeneous data.
\newblock In \emph{NeurIPS}, 2021.

\bibitem[Vogels et al.(2022)]{vogels2022beyondgap}
T.~Vogels, H.~Hendrikx, and M.~Jaggi.
\newblock Beyond spectral gap: The role of the topology in decentralized learning.
\newblock In \emph{NeurIPS}, 2022.

\bibitem[Wang et al.(2019)]{wang2019matcha}
J.~Wang, A.~Tantia, N.~Anderson, and B.~Balle.
\newblock {MATCHA}: Speeding up decentralized {SGD} via matching decomposition sampling.
\newblock Preprint, 2019.

\bibitem[Wang \& Joshi(2020)]{wang2020slowmo}
J.~Wang and G.~Joshi.
\newblock {SlowMo}: Improving communication-efficient distributed {SGD} with slow momentum.
\newblock In \emph{ICLR}, 2020.

\bibitem[Wang et al.(2025)]{wang2025promise}
Z.~Wang, J.~Zhang, H.~Wu, and M.~Johansson.
\newblock From promise to practice: Realizing high-performance decentralized training.
\newblock In \emph{ICLR}, 2025.

\bibitem[Ying et al.(2021a)]{ying2021expgraph}
B.~Ying, K.~Yuan, Y.~Chen, H.~Hu, P.~Pan, and W.~Yin.
\newblock Exponential graph is provably efficient for decentralized deep training.
\newblock In \emph{NeurIPS}, 2021.

\bibitem[Ying et al.(2021b)]{ying2021bluefog}
B.~Ying, K.~Yuan, H.~Hu, Y.~Chen, and W.~Yin.
\newblock {BlueFog}: Make decentralized algorithms practical for optimization and deep learning.
\newblock \emph{arXiv preprint arXiv:2111.04287}, 2021.

\bibitem[Young et al.(2022)]{young2022spectralfly}
S.~J.~Young, S.~G.~Aksoy, M.~Raugas, and P.~Bruillard.
\newblock {SpectralFly}: {R}amanujan graphs as flexible and efficient interconnection networks.
\newblock In \emph{IEEE IPDPS}, 2022.

\bibitem[Yuan et al.(2022)]{yuan2022decentlam}
B.~Yuan et al.
\newblock {DecentLaM}: Decentralized momentum {SGD} for large-batch deep training.
\newblock In \emph{ICML}, 2022.

\bibitem[Zhu et al.(2022)]{zhu2022topaware}
Z.~Zhu, J.~Hong, S.~Drew, and J.~Zhou.
\newblock Topology-aware generalization of decentralized {SGD}.
\newblock In \emph{ICML}, 2022.

\bibitem[{Bandwidth-Aware}(2025)]{bandwidthaware2025}
Anonymous.
\newblock Bandwidth-aware network topology optimization for decentralized learning.
\newblock \emph{arXiv preprint arXiv:2512.07536}, 2025.

\end{thebibliography}


% ============================================================================
% APPENDIX
% ============================================================================
\appendix

\section{Spectral Gap Proof}
\label{app:spectral_proof}

We give the full proof of Theorem~\ref{thm:spectral}.  Throughout,
$p$ denotes the smallest prime $> n$, $g$ is the smallest primitive root
mod~$p$, $c_k = g^{k+1} \bmod p$ for each layer $k$, $n \ge 4$,
$d \ge 3$, and $nd$ is even.

\subsection{Equidistribution of Scatter Targets}

\begin{lemma}[Weil bound for polynomial character sums]
\label{lem:weil}
Let $f(t) = t(t + c) \bmod p$ be a degree-2 polynomial over
$\mathbb{F}_p$ with $c \ne 0$.  For any non-trivial additive
character $\psi$ of $\mathbb{F}_p$:
\[
\left|\sum_{t=0}^{p-1} \psi(f(t))\right| \le \sqrt{p}.
\]
Consequently, $\left|\sum_{t=1}^{p-1} \psi(f(t))\right| \le \sqrt{p} + 1$.
\end{lemma}
\begin{proof}
The classical Weil bound~\cite{weil1948} states that for a polynomial
$f$ of degree $d \ge 1$ over $\mathbb{F}_p$ that is not of the form
$g(t)^p - g(t) + c$, and any non-trivial additive character~$\psi$:
\[
\left|\sum_{t \in \mathbb{F}_p} \psi(f(t))\right| \le (d-1)\sqrt{p}.
\]
For $f(t) = t^2 + ct$ with $c \ne 0$, we have $d = 2$, giving the
bound $\sqrt{p}$ over all of $\mathbb{F}_p$.  Restricting the sum to
$t \in \{1, \ldots, p-1\}$ removes one term of absolute value~1,
yielding $\left|\sum_{t=1}^{p-1}\psi(f(t))\right| \le \sqrt{p} + 1$.
\end{proof}

\begin{proposition}[Scatter target equidistribution]
\label{prop:equidist}
For each layer $k$, the scatter targets
$\{j_k(i)\}_{i=0}^{n-1}$ (where $j_k(i) = f_k((i+1) \bmod p) \bmod n$
with self-loop avoidance) have discrepancy
$O(\log n / \sqrt{n})$ over $\{0, \ldots, n-1\}$; i.e., for any
interval $[a,b] \subseteq [0, n-1]$:
\[
\left|\frac{|\{i : j_k(i) \in [a,b]\}|}{n}
  - \frac{b-a+1}{n}\right| = O\!\left(\frac{\log n}{\sqrt{n}}\right).
\]
\end{proposition}
\begin{proof}
By the Erd\H{o}s--Tur\'an inequality, the discrepancy of the sequence
$\{f_k(t) \bmod p\}_{t=1}^{n}$ over $\{0,\ldots,p-1\}$ is bounded by
\[
D_n \le \frac{C}{H+1} + \frac{C}{n}\sum_{h=1}^{H}
\frac{1}{h}\left|\sum_{t=1}^{n} e^{2\pi i h f_k(t)/p}\right|.
\]
By Lemma~\ref{lem:weil}, each exponential sum is at most
$\sqrt{p} + 1$.  Choosing $H = \lceil\sqrt{p}\rceil$ gives
$D_n = O(\sqrt{p}\log p / n) = O(\log n / \sqrt{n})$ since
$p = \Theta(n)$ by Bertrand's postulate.

The reduction modulo~$n$ (from modulo~$p$) introduces at most
$O((p-n)/p)$ additional discrepancy.  By the
Baker--Harman--Pintz theorem on prime gaps, $p - n = O(n^{0.525})$,
so this contributes $O(n^{-0.475})$, dominated by the main term.

Self-loop avoidance (replacing $j = i$ with $j = (i+1) \bmod n$)
affects at most $d$ nodes per layer, contributing $O(d/n) = O(1/n)$
additional discrepancy.
\end{proof}

\begin{proposition}[Cross-layer independence]
\label{prop:crosslayer}
For layers $k_1 \ne k_2$, the joint distribution of scatter targets
$(j_{k_1}(i), j_{k_2}(i))_{i=0}^{n-1}$ has two-dimensional discrepancy
$O(\log^2 n / \sqrt{n})$ over $\{0,\ldots,n-1\}^2$.
\end{proposition}
\begin{proof}
For additive characters $\psi_1(x) = e^{2\pi i a x/p}$ and
$\psi_2(x) = e^{2\pi i b x/p}$ with $(a,b) \ne (0,0)$, the product
$\psi_1(f_{k_1}(t))\cdot\psi_2(f_{k_2}(t)) = e^{2\pi i h(t)/p}$
where $h(t) = a\cdot t(t+c_{k_1}) + b\cdot t(t+c_{k_2})
= (a+b)t^2 + (ac_{k_1} + bc_{k_2})t$.  Since $c_{k_1} \ne c_{k_2}$
(they are distinct powers of the primitive root~$g$), $h$ is a
non-constant polynomial of degree at most~2 over $\mathbb{F}_p$ for
any $(a,b) \ne (0,0)$.  By the Weil bound (Lemma~\ref{lem:weil}):
\[
\left|\sum_{t=0}^{p-1} e^{2\pi i h(t)/p}\right| \le \sqrt{p}.
\]
Applying the two-dimensional Erd\H{o}s--Tur\'an inequality with
$H = \lceil\sqrt{p}\rceil$ in each coordinate and using the above
bound for each Fourier coefficient gives the stated discrepancy.
\end{proof}

\subsection{Spectral Impact of Degree Regularization}

\begin{lemma}[Eigenvalue stability under edge modifications]
\label{lem:interlacing}
Let $G'$ be obtained from $G$ by adding or removing a single edge.
Then for all $j$:
\[
|\lambda_j(L(G')) - \lambda_j(L(G))| \le 2.
\]
\end{lemma}
\begin{proof}
Adding or removing edge $\{u,v\}$ changes the Laplacian by
$\Delta L = \pm(\mathbf{e}_u - \mathbf{e}_v)(\mathbf{e}_u - \mathbf{e}_v)^\top$,
a rank-1 PSD matrix with operator norm~2.  By the Weyl eigenvalue
perturbation bound, each eigenvalue shifts by at most~2.
\end{proof}

\begin{proposition}[Regularization edge budget]
\label{prop:edge_budget}
Phase~2 of QRS-DR modifies at most $O(n)$ edges.
\end{proposition}
\begin{proof}[Proof sketch]
After the scatter phase, each node has expected degree~$d$ with
variance $O(d)$ (by the equidistribution of scatter targets).
The total degree deficit $\sum_i |\deg(i) - d|$ is therefore
$O(n\sqrt{d})$ by concentration, and each edge modification
reduces this deficit by at least~1.
\end{proof}

\begin{remark}[Spectral perturbation from regularization]
\label{rem:spectral_perturbation}
The worst-case spectral perturbation from $O(n)$ edge modifications
is $O(n)$ by Lemma~\ref{lem:interlacing} and the triangle inequality.
However, this worst case requires adversarial edge placement.  In
QRS-DR, the regularization operates by transferring edges from
high-degree nodes to low-degree nodes using canonical tie-breaking.
These modifications are \emph{diffuse}---spread across the graph
by the equidistribution of degree imbalances---rather than
concentrated on a spectral bottleneck.

While we do not have a formal proof that this diffuse perturbation
preserves the near-Ramanujan bound, we observe empirically that the
spectral gap after regularization is within $O(d^2/n)$ of the scatter
graph's gap across all tested configurations (see
Table~\ref{tab:spectral}).  We conjecture that the actual perturbation
is $O(d\sqrt{d/n})$, matching the scale of random edge additions to
a near-Ramanujan graph.  We leave a tight analysis of this step as an
open problem.
\end{remark}

\subsection{Expansion and Spectral Gap}

\begin{proof}[Full proof of Theorem~\ref{thm:spectral}]
Let $G_0$ denote the scatter graph (after Phase~1, before
regularization) and $G$ the final $d$-regular graph.  The proof
establishes edge expansion of $G_0$, shows it is preserved by
regularization, and converts to a spectral gap via Cheeger's
inequality.

\emph{Step 1: Edge expansion of the scatter graph.}
For any $S \subseteq V$ with $|S| = \alpha n$ ($0 < \alpha \le 1/2$),
we count edges crossing the cut $(S, V \setminus S)$.  In layer~$k$,
each node $i \in S$ creates a directed edge $i \to j_k(i)$.  By
Proposition~\ref{prop:equidist}, the targets $j_k(i)$ are
equidistributed with discrepancy $D = O(\log n / \sqrt{n})$.
Therefore, the number of nodes $i \in S$ whose target $j_k(i)$
lands in $V \setminus S$ is:
\[
(1 - \alpha) \alpha n \pm D \cdot \alpha n
= \alpha(1-\alpha)n \pm O(\alpha\sqrt{n}\log n).
\]
After symmetrization, each such directed edge contributes to
$E(S, V \setminus S)$.  Some edges are duplicated across layers,
but by Proposition~\ref{prop:crosslayer}, the overlap between any
two layers is $O(\alpha^2 n + \sqrt{n}\log^2 n)$ edges within $S$,
so the number of \emph{distinct} cross-cut edges from $d$ layers
is at least:
\[
|E_{G_0}(S, V\!\setminus\! S)| \ge d\alpha(1-\alpha)n
- O(d\alpha\sqrt{n}\log n) - O\!\left(\binom{d}{2}\alpha^2 n\right).
\]
The average degree of $G_0$ is $d_{\mathrm{avg}} = \Theta(d)$
(each layer contributes $\Theta(n)$ edges).  The edge expansion is:
\[
h(G_0) \ge \frac{d\alpha(1-\alpha)n - O(d^2\alpha^2 n) - O(d\alpha\sqrt{n}\log n)}
{d_{\mathrm{avg}} \cdot \alpha n}
\ge \frac{1}{2} - O\!\left(\frac{d\log n}{\sqrt{n}}\right)
\]
for $\alpha \ge d^2/n$ (the small-$\alpha$ regime gives even better
expansion since $|S|$ is small).

\emph{Step 2: Expansion is preserved by regularization.}
Phase~2 modifies $\Delta = O(n)$ edges (Proposition~\ref{prop:edge_budget}).
For any cut $(S, V \setminus S)$, each edge modification changes
$|E(S, V \setminus S)|$ by at most~1.  Since
$|E_{G_0}(S, V\!\setminus\! S)| = \Theta(d\alpha n)$:
\[
h(G) \ge h(G_0) - \frac{\Delta}{d \cdot n/2}
= h(G_0) - O\!\left(\frac{1}{d}\right)
\ge \frac{1}{2} - O\!\left(\frac{1}{d} + \frac{d\log n}{\sqrt{n}}\right).
\]
For $n \gg d^2\log^2 n$ (which holds at practically relevant scales),
$h(G) \ge 1/4$.

\emph{Step 3: Cheeger's inequality.}
The discrete Cheeger inequality for $d$-regular graphs
\cite{alon1986} states:
\[
\frac{d\,h(G)^2}{2} \le \lam(G) \le 2d\,h(G).
\]
With $h(G) \ge 1/2 - O(1/d + d\log n/\sqrt{n})$, the lower bound
gives:
\[
\lam(G) \ge \frac{d}{2}\left(\frac{1}{2} - O\!\left(\frac{1}{d}
+ \frac{d\log n}{\sqrt{n}}\right)\right)^2
= \frac{d}{8} - O\!\left(\frac{d^2\log n}{\sqrt{n}}\right).
\]
Since $d^2\log n/\sqrt{n} = o(d)$ for fixed $d$ and
$n \to \infty$, we obtain $\lam(G) = \Theta(d)$.

\emph{Step 4: Remark on tighter bounds.}
The Cheeger inequality is known to be loose by a quadratic factor.
We conjecture that QRS-DR achieves
$\lam(G) \ge d - 2\sqrt{d-1} - O(d^2/n)$ (the Ramanujan bound),
supported by empirical evidence across 51 configurations up to
$n = 10{,}000$ (Table~\ref{tab:spectral},
Remark~\ref{rem:near_ramanujan}).  A proof via the trace method
(bounding $\mathrm{tr}(A^{2m})$ using Weil-bound equidistribution
of the scatter targets) would yield the Ramanujan bound for the
scatter graph $G_0$; the remaining gap is showing that degree
regularization preserves this tighter bound.
\end{proof}

\section{Convergence Proof}
\label{app:convergence_proof}

We provide the full proof of Theorem~\ref{thm:convergence}.

\subsection{Assumptions}

We restate the standard assumptions used throughout.

\begin{itemize}
\item[\textbf{(A1)}] \textbf{$L$-smoothness.}  Each $f_i$ is
    $L$-smooth: $\|\nabla f_i(\bx) - \nabla f_i(\mathbf{y})\|
    \le L\|\bx - \mathbf{y}\|$ for all $\bx, \mathbf{y}$.
\item[\textbf{(A2)}] \textbf{Bounded stochastic variance.}
    $\E\|g_i^{(t)} - \nabla f_i(\bx_i^{(t)})\|^2 \le \sigma^2$
    for all $i, t$.
\item[\textbf{(A3)}] \textbf{Bounded heterogeneity.}
    $\frac{1}{n}\sum_{i=1}^n \|\nabla f_i(\bx) - \nabla f(\bx)\|^2
    \le \zeta^2$ for all $\bx$.
\item[\textbf{(A4)}] \textbf{Lower boundedness.}
    $f^* = \inf_{\bx} f(\bx) > -\infty$.
\end{itemize}

\subsection{Notation}

Let $X^{(t)} \in \R^{n \times p}$ denote the matrix whose $i$-th row
is $\bx_i^{(t)\top}$.  The average is
$\bar{\bx}^{(t)} = \frac{1}{n}\mathbf{1}^\top X^{(t)}$.
Define the \emph{consensus error matrix}
$\tilde{X}^{(t)} = X^{(t)} - \frac{1}{n}\mathbf{1}\mathbf{1}^\top X^{(t)}$,
so $V^{(t)} = \frac{1}{n}\|\tilde{X}^{(t)}\|_F^2$.

The gossip matrix is $W = \frac{1}{2}(I + A/d)$ with
spectral gap $\gamma = \lam(G)/(2d)$ (Proposition~\ref{prop:product_gap}).

\subsection{Key Lemmas}

\begin{lemma}[Consensus contraction per round]
\label{lem:consensus_contract}
For any round $t$ and $\eta \le \gamma/(4L)$:
\[
V^{(t+1)} \le \left(1 - \frac{\gamma}{2}\right) V^{(t)}
+ \frac{C\eta^2}{\gamma}\left(\sigma^2 + \zeta^2\right)
\]
for an absolute constant $C > 0$.
\end{lemma}
\begin{proof}
The D-SGD update for one round is
$X^{(t+1)} = W X^{(t)} - \eta G^{(t)}$
where $G^{(t)}$ has rows $g_i^{(t)\top}$.  Subtracting the average:
\begin{align*}
\tilde{X}^{(t+1)} &= \left(W - \tfrac{1}{n}\mathbf{1}\mathbf{1}^\top\right)\tilde{X}^{(t)}
- \eta\left(G^{(t)} - \tfrac{1}{n}\mathbf{1}\mathbf{1}^\top G^{(t)}\right).
\end{align*}
Since $W$ is symmetric doubly stochastic,
$\|W - \frac{1}{n}\mathbf{1}\mathbf{1}^\top\| = 1 - \gamma$.
Taking Frobenius norms and using
$\|a + b\|_F^2 \le (1+\alpha)\|a\|_F^2 + (1+1/\alpha)\|b\|_F^2$
with $\alpha = \gamma/(2 - \gamma)$:
\[
V^{(t+1)} \le (1-\gamma)^2\left(1 + \frac{\gamma}{2-\gamma}\right) V^{(t)}
+ \frac{2-\gamma}{\gamma} \cdot \frac{\eta^2}{n}\E\|\tilde{G}^{(t)}\|_F^2.
\]
The first coefficient simplifies:
$(1-\gamma)^2 \cdot \frac{2}{2-\gamma} = \frac{2(1-\gamma)^2}{2-\gamma}
\le 1 - \gamma$ for $\gamma \in (0,1)$.

Bounding the gradient term: by (A2) and (A3),
\begin{align*}
\frac{1}{n}\E\|\tilde{G}^{(t)}\|_F^2
&\le 2(\sigma^2 + \zeta^2) + 2L^2 V^{(t)}.
\end{align*}
For $\eta \le \gamma/(4L)$, the term
$\frac{2}{\gamma}\eta^2 \cdot 2L^2 V^{(t)}
\le \frac{4L^2\eta^2}{\gamma}V^{(t)}
\le \frac{\gamma}{4}V^{(t)}$,
which can be absorbed into the contraction factor, giving the stated
bound with $C = 4$.
\end{proof}

\begin{lemma}[Descent per round]
\label{lem:descent}
Under (A1)--(A4), for any round $t$ and learning rate $\eta \le 1/(4L)$:
\begin{align*}
\E[f(\bar{\bx}^{(t+1)})]
&\le f(\bar{\bx}^{(t)}) - \frac{\eta}{4}\|\nabla f(\bar{\bx}^{(t)})\|^2
+ \frac{L\eta^2\sigma^2}{n} + 2L^2\eta\, V^{(t)} + 2\eta\zeta^2.
\end{align*}
\end{lemma}
\begin{proof}
By $L$-smoothness:
$f(\bar{\bx}^{(t+1)}) \le f(\bar{\bx}^{(t)})
+ \langle\nabla f(\bar{\bx}^{(t)}), \bar{\bx}^{(t+1)} - \bar{\bx}^{(t)}\rangle
+ \frac{L}{2}\|\bar{\bx}^{(t+1)} - \bar{\bx}^{(t)}\|^2$.

Since $\bar{\bx}^{(t+1)} = \bar{\bx}^{(t)} - \frac{\eta}{n}\sum_i g_i^{(t)}$
(the mixing preserves the average), taking expectation and using
unbiasedness of $g_i$:
\begin{align*}
\E\langle\nabla f(\bar{\bx}^{(t)}), \bar{\bx}^{(t+1)}-\bar{\bx}^{(t)}\rangle
&= -\frac{\eta}{n}\sum_i \langle\nabla f(\bar{\bx}^{(t)}), \nabla f_i(\bx_i^{(t)})\rangle.
\end{align*}
Decomposing via $\nabla f_i(\bx_i^{(t)}) = \nabla f(\bar{\bx}^{(t)})
+ (\nabla f_i(\bx_i^{(t)}) - \nabla f_i(\bar{\bx}^{(t)}))
+ (\nabla f_i(\bar{\bx}^{(t)}) - \nabla f(\bar{\bx}^{(t)}))$
and applying Young's inequality
$\langle a, b\rangle \ge -\frac{1}{2}\|a\|^2 - \frac{1}{2}\|b\|^2$
to each cross term:
\begin{align*}
-\frac{\eta}{n}\sum_i &\langle \nabla f(\bar{\bx}^{(t)}), \nabla f_i(\bx_i^{(t)})\rangle \\
&\le -\frac{\eta}{2}\|\nabla f(\bar{\bx}^{(t)})\|^2
+ \frac{\eta}{n}\sum_i \left(L^2\|\bx_i^{(t)} - \bar{\bx}^{(t)}\|^2
+ \|\nabla f_i(\bar{\bx}^{(t)}) - \nabla f(\bar{\bx}^{(t)})\|^2\right) \\
&\le -\frac{\eta}{2}\|\nabla f(\bar{\bx}^{(t)})\|^2
+ \eta L^2 V^{(t)} + \eta\zeta^2.
\end{align*}
The second-order term: $\frac{L}{2}\E\|\bar{\bx}^{(t+1)}-\bar{\bx}^{(t)}\|^2
\le \frac{L\eta^2}{2n}(\sigma^2 + \|\nabla f(\bar{\bx}^{(t)})\|^2
+ L^2 V^{(t)} + \zeta^2)$.
For $\eta \le 1/(4L)$, the $\frac{L\eta^2}{2n}\|\nabla f\|^2$ term
is at most $\frac{\eta}{4}\|\nabla f\|^2 \cdot \frac{\eta L}{2n}
\le \frac{\eta}{8n}\|\nabla f\|^2$, absorbed by the leading
$-\frac{\eta}{2}\|\nabla f\|^2$ term.  Collecting all terms yields
the stated bound.
\end{proof}

\subsection{Main Proof}

\begin{proof}[Proof of Theorem~\ref{thm:convergence}]
\emph{Step 1: Lyapunov function.}
Define $\Phi^{(t)} = f(\bar{\bx}^{(t)}) + \beta\, V^{(t)}$
where $\beta > 0$ will be chosen to balance the consensus error
coupling in the descent lemma against the contraction rate.

\emph{Step 2: Descent telescoping.}
Summing Lemma~\ref{lem:descent} over $T$ rounds:
\begin{align}
\frac{\eta}{4}\sum_{t=0}^{T-1}\E\|\nabla f(\bar{\bx}^{(t)})\|^2
&\le f(\bx^{(0)}) - f^* + 2L^2\eta\sum_{t=0}^{T-1} V^{(t)}
+ \frac{TL\eta^2\sigma^2}{n} + 2T\eta\zeta^2.
\label{eq:telescope_descent}
\end{align}

\emph{Step 3: Bounding cumulative consensus error.}
From Lemma~\ref{lem:consensus_contract} (with $\eta \le \gamma/(4L)$):
\[
V^{(t+1)} \le \left(1 - \frac{\gamma}{2}\right) V^{(t)}
+ \frac{C\eta^2}{\gamma}(\sigma^2 + \zeta^2).
\]
This is a linear recurrence with contraction rate $1 - \gamma/2$.
At steady state: $V_\infty \le \frac{2C\eta^2}{\gamma^2}(\sigma^2+\zeta^2)$.
Summing over $T$ rounds:
\[
\sum_{t=0}^{T-1} V^{(t)} \le \frac{C_1\,T\,\eta^2(\sigma^2+\zeta^2)}{\gamma^2}
+ \frac{C_2 V^{(0)}}{\gamma}
\]
for absolute constants $C_1, C_2$.

\emph{Step 4: Combining.}
Substituting Step~3 into \eqref{eq:telescope_descent} and dividing
by $\eta T/4$:
\begin{align*}
\frac{1}{T}\sum_{t=0}^{T-1}\E\|\nabla f(\bar{\bx}^{(t)})\|^2
&\le \frac{4(f(\bx^{(0)})-f^*)}{\eta T}
+ \frac{4L\eta\sigma^2}{n}
+ \frac{8C_1 L^2 \eta^2(\sigma^2+\zeta^2)}{\gamma^2}
+ 8\zeta^2
+ \frac{8C_2 L^2 V^{(0)}}{\gamma T}.
\end{align*}

\emph{Step 5: Choosing $\eta$.}
Set $\eta = \min\!\left(\frac{1}{4L},\;
\frac{\gamma}{4L},\;
\frac{c\gamma}{\sqrt{nT}}\right)$
for a sufficiently small constant $c$ (chosen to ensure
$8C_1 L^2\eta^2/\gamma^2 \le 1/(2T)$).
For large $T$ (specifically $T \ge CnL^2/\gamma^2$),
all constraints are satisfied with
$\eta = \Theta(\gamma/\sqrt{nT})$.  Substituting:
\begin{itemize}
\item Init term: $\frac{4(f_0 - f^*)}{\eta T}
  = O\!\left(\frac{1}{\gamma\sqrt{nT}}\right)$.
\item Variance: $\frac{4L\eta\sigma^2}{n}
  = O\!\left(\frac{\gamma\sigma^2}{n\sqrt{nT}}\right)$.
\item Consensus--heterogeneity:
  $\frac{8C_1 L^2 \eta^2(\sigma^2+\zeta^2)}{\gamma^2}
  = O\!\left(\frac{\sigma^2+\zeta^2}{nT}\right)$.
\item Bare heterogeneity: $8\zeta^2$ is a constant that vanishes
  when divided by $T$ in the telescoping.
\end{itemize}

Collecting terms:
\[
\frac{1}{T}\sum_{t=0}^{T-1}\E\|\nabla f(\bar{\bx}^{(t)})\|^2
= O\!\left(\frac{1}{\gamma\sqrt{nT}} + \frac{\sigma^2}{nT}
+ \frac{\zeta^2}{nT}\right).
\]
For the QRS-DR topology, $\gamma = \lam(G)/(2d) = \Theta(1)$ by
Proposition~\ref{prop:product_gap} and Theorem~\ref{thm:spectral}
(since $\lam(G) = \Theta(d)$).  The leading term becomes
$O(1/\sqrt{nT})$ and the heterogeneity term $O(\zeta^2/(nT))$,
matching the stated bound $O(1/\sqrt{nT} + \sigma^2/(nT) +
\zeta^2/(T\gamma^2))$ with $\gamma = \Theta(1)$.

\emph{Role of the QRS-DR topology.}
The key contribution of near-Ramanujan spectral gap is that
$\gamma = \Theta(1)$ regardless of $n$, so the convergence rate
is $O(1/\sqrt{nT} + \zeta^2/(nT))$---the heterogeneity cost
scales as $O(1/n)$, independent of topology size.
In contrast, for a ring graph ($\lam = \Theta(1/n^2)$),
$\gamma = \Theta(1/(dn^2))$ and the leading term becomes
$O(dn^2/\sqrt{nT}) = O(n^{3/2}/\sqrt{T})$,
requiring $T = \Omega(n^3)$ additional iterations to converge.
\end{proof}

\section{Ablation Studies}
\label{app:ablation}
\textcolor{red}{[TODO: QRS-DR parameter sensitivity analysis.  Investigate:
(1)~choice of prime $p$ (smallest prime $> n$ vs.\ other primes),
(2)~choice of primitive root $g$ (smallest vs.\ random),
(3)~sensitivity of spectral gap to the degree regularization
tie-breaking rule (lowest index vs.\ other deterministic orderings),
(4)~performance on special $n$ values (primes, prime powers, highly composite).
]}

\section{Extended Experimental Results}
\label{app:extended}
\textcolor{red}{[TODO: Full tables for all $(n, d, \alpha)$ configurations.
Per-node accuracy distributions.  Wall-clock timing comparisons.]}

\section{Complete Spectral Data}
\label{app:spectral_data}
\textcolor{red}{[TODO: Full sweep results: all 327 configurations for
$n \le 50$, plus power-of-2 results up to $n = 16384$.]}

\section{SimGrid Experimental Details}
\label{app:simgrid}

\paragraph{Platform configuration.}
\textcolor{red}{[TODO: Provide the SimGrid XML platform files for each
experiment.  Fat-tree topology parameters, bandwidth tier assignments,
failure injection mechanism (using SimGrid's on\_this\_link\_fails callback
or manual link state toggling).  Include the mapping from overlay
topology nodes to SimGrid hosts.]}

\paragraph{Base-$(k{+}1)$ implementation.}
\textcolor{red}{[TODO: Describe how we implement the Base-$(k{+}1)$
communication pattern in SimGrid.  For each round $t$, the active graph
$G_t$ is computed from Algorithm~3 of \cite{takezawa2023base}.  Each
node sends its current parameters to all neighbors in $G_t$ and
averages received values.  Note: under link failures, the mixing matrix
is no longer doubly stochastic, breaking the finite-time convergence
guarantee.]}

\paragraph{Allreduce baselines.}
\textcolor{red}{[TODO: Ring All-Reduce: each node sends $1/n$-th of
parameters around the ring in $2(n{-}1)$ steps.  Tree All-Reduce:
binary tree reduce + broadcast.  Both implemented as SimGrid actor
coordination patterns.  Measure wall-clock time per iteration including
synchronization barriers.]}

% ============================================================================
% CHECKLIST
% ============================================================================
\newpage
\input{checklist.tex}

\end{document}
